% Unit 8 (Acids and Bases) guided notes -- 9 blocks, CED sequence.
\documentclass{apchem-notes}

\apcunit{8}
\apcunittitle{Acids and Bases}
\apcdoctitle{Guided Notes}
\apcsubtitle{CED 8.1--8.11 \textbullet\ Zumdahl Ch 14--15, \S16.1 \textbullet\ 9 blocks}

\begin{document}
\apctitleblock

\begin{apcnote}[The question to ask before every calculation]
This unit contains four situations, and nearly every lost mark comes from
running the wrong one. Before touching a calculator, name the major species
in solution and decide which case you are in:

\begin{center}
\renewcommand{\arraystretch}{1.2}
\begin{tabular}{@{}lll@{}}
\hline
\textbf{What is in solution} & \textbf{Case} & \textbf{Method} \\
\hline
strong acid or base alone & complete ionization & pH directly \\
weak acid (or base) alone & one equilibrium & ICE with $K_a$ (or $K_b$) \\
weak acid \emph{and} its conjugate & buffer &
  Henderson--Hasselbalch \\
strong reagent added & react to completion \emph{first} & then one of the
  above \\
\hline
\end{tabular}
\end{center}

At \textbf{11--15\%} this is the second-heaviest unit on the exam. It is
also the one where method selection, rather than arithmetic, decides the
grade.
\end{apcnote}

% =====================================================================
\blocklesson{Br\o nsted--Lowry and Conjugate Pairs}{CED 8.1 \textbullet\ Zumdahl \S14.1--14.2}

\begin{objectives}
  \item Identify conjugate acid--base pairs.
  \item Relate acid strength to $K_a$ and to conjugate base strength.
\end{objectives}

\reading{Zumdahl \S14.1--14.2}{697--705}

\begin{warmup}[4]
  \item Net ionic equation for strong acid $+$ strong base:
        \answerblank[4.6cm]{\ce{H+ + OH^- -> H2O}}
  \item Name two strong acids (there are six):
        \answerblank[4cm]{\ce{HCl}, \ce{HNO3}}
  \item $K$ expression for \ce{HA <=> H+ + A^-}:
        \answerblank[4.2cm]{$[\ce{H+}][\ce{A^-}]/[\ce{HA}]$}
  \item $K \ll 1$ favors: \answerblank[2.4cm]{reactants}
\end{warmup}

\segment{INSTRUCTION A}{25}{Two models, one of which AP uses}

\notesection{Proton donors and acceptors}{8.1}
\practice{1}

\begin{itemize}[leftmargin=*,itemsep=0.3em]
  \item \textbf{Arrhenius:} an acid produces \answerblank[2cm]{\ce{H+}} in
        water; a base produces \answerblank[2cm]{\ce{OH-}}. Historically
        important, but it applies only to aqueous solutions.
  \item \textbf{Br\o nsted--Lowry:} an acid is a proton
        \answerblank[2.2cm]{donor}; a base is a proton
        \answerblank[2.4cm]{acceptor}. This is the model the CED uses.
\end{itemize}

Water acts as a base because its oxygen carries
\answerblank[5cm]{unshared electron pairs}, either of which can bond to a
proton. So \ce{H+} in solution is really \term{hydronium}, \ce{H3O+};
writing \ce{H+} is accepted shorthand.

\notesection{Conjugate pairs}{8.1}

Every Br\o nsted--Lowry reaction contains \emph{two} pairs differing by one
proton:
\[ \ce{HC2H3O2 + H2O <=> H3O+ + C2H3O2^-} \]
\ce{HC2H3O2}/\ce{C2H3O2^-} is one pair; \ce{H2O}/\ce{H3O+} is the other.

\begin{apctrap}
\textbf{The stronger the acid, the weaker its conjugate base.} \ce{HCl} is
strong, so \ce{Cl-} is such a poor base it is effectively inert in water.
\ce{HCN} is very weak, so \ce{CN-} is a base strong enough to matter.

Quantitatively, $K_a \times K_b$ equals \answerblank[2cm]{$K_w$} for a
conjugate pair --- adding the two reactions gives the autoionization of
water.
\end{apctrap}

\segment{GUIDED PRACTICE}{15}{Identifying pairs}

\begin{enumerate}[leftmargin=*,itemsep=0.3em]
  \item Conjugate base of \ce{HNO2}: \answerblank[2.6cm]{\ce{NO2^-}}
  \item Conjugate acid of \ce{NH3}: \answerblank[2.6cm]{\ce{NH4+}}
  \item Conjugate base of \ce{H2PO4^-}:
        \answerblank[2.6cm]{\ce{HPO4^2-}}
  \item Conjugate acid of \ce{H2O}: \answerblank[2.6cm]{\ce{H3O+}}
  \item Rank as bases, weakest first: \ce{Cl-}, \ce{F-}, \ce{CN-}
        \answerblank[5cm]{\ce{Cl-} $<$ \ce{F-} $<$ \ce{CN-}}
\end{enumerate}

\segment{APPLICATION}{20}{Strength and $K_a$}

\begin{enumerate}[leftmargin=*,itemsep=0.4em]
  \item A larger $K_a$ means what about the acid, and about its conjugate
        base? \answerlines[2]{A larger $K_a$ means a stronger acid --- more
        of it ionizes. Its conjugate base is correspondingly weaker,
        because an anion that gave up its proton readily has little
        affinity for taking it back.}
  \item Given $K_a(\ce{HC2H3O2}) = 1.8\times10^{-5}$, find $K_b$ for
        acetate. \workspace[1.8cm]
        \keyonly{\begin{apcanswer}$K_b = K_w/K_a =
        1.0\times10^{-14}/1.8\times10^{-5} =
        \mathbf{5.6\times10^{-10}}$\end{apcanswer}}
\end{enumerate}

\begin{exitticket}
Explain why \ce{Cl-} does not act as a base in water but \ce{F-} does.
\keyonly{\begin{apcanswer}\ce{HCl} is a strong acid, so \ce{Cl-} has
essentially no affinity for protons and does not remove them from water.
\ce{HF} is weak, so \ce{F-} readily accepts a proton from water, generating
\ce{OH-} and making the solution basic.\end{apcanswer}}
\end{exitticket}

% =====================================================================
\blocklesson{pH, pOH, and Strong Acids and Bases}{CED 8.2 \textbullet\ Zumdahl \S14.3--14.4, 14.6}

\begin{objectives}
  \item Convert among $[\ce{H+}]$, $[\ce{OH-}]$, pH and pOH.
  \item Calculate the pH of strong acid and strong base solutions.
\end{objectives}

\reading{Zumdahl \S14.3--14.4, 14.6}{705--724}

\begin{warmup}[3]
  \item $K_aK_b = $ \answerblank[1.8cm]{$K_w$}
  \item Conjugate base of \ce{HF}: \answerblank[2cm]{\ce{F-}}
  \item $K_w$ at \SI{25}{\celsius}:
        \answerblank[3cm]{$1.0\times10^{-14}$}
\end{warmup}

\segment{INSTRUCTION A}{25}{The logarithmic scale}

\notesection{pH, pOH, and $K_w$}{8.2}
\practice{5}

\[ \mathrm{pH} = -\log[\ce{H+}] \qquad
   \mathrm{pOH} = -\log[\ce{OH-}] \qquad
   \mathrm{pH} + \mathrm{pOH} = 14.00 \]

Each pH unit is a factor of \answerblank[1.6cm]{10} in $[\ce{H+}]$, so
pH 3 is \answerblank[2cm]{1000} times more acidic than pH 6.

\begin{apctrap}
\textbf{Neutral means $[\ce{H+}] = [\ce{OH-}]$, not pH 7.} $K_w$ rises with
temperature, so pure water at \SI{50}{\celsius} has a pH below 7 and is
still perfectly neutral. The number 7 applies only at \SI{25}{\celsius}.
\end{apctrap}

\segment{INSTRUCTION B}{20}{When ionization is complete}

\notesection{Strong acids and bases}{8.2}
\practice{5}

For a strong acid, ionization is complete, so $[\ce{H+}]$ simply equals the
\answerblank[4.2cm]{acid concentration}. No equilibrium calculation is
needed.

There are only \textbf{six strong acids}, worth memorizing outright:
\[ \ce{HCl} \quad \ce{HBr} \quad \ce{HI} \quad \ce{HNO3} \quad
   \ce{H2SO4} \quad \ce{HClO4} \]
Note that \ce{HF} is \emph{not} among them. The strong bases are the Group 1
hydroxides and \ce{Ca(OH)2}, \ce{Sr(OH)2}, \ce{Ba(OH)2}.

\begin{workedexample}[Worked example 1: the factor everyone forgets]
\SI{0.010}{\Molar} \ce{HCl}: $[\ce{H+}] = 0.010$, so
$\mathrm{pH} = \mathbf{2.00}$.

\SI{0.010}{\Molar} \ce{Ba(OH)2} produces \emph{two} hydroxides per formula
unit:
\[ [\ce{OH-}] = 2(0.010) = 0.020 \quad\Longrightarrow\quad
   \mathrm{pOH} = 1.70,\ \mathrm{pH} = \mathbf{12.30} \]
--- not 12.00. Always check the formula for how many ions are released.
\end{workedexample}

\segment{APPLICATION}{20}{Strong acid and base calculations}

\begin{enumerate}[leftmargin=*,itemsep=0.4em]
  \item pH of \SI{0.0025}{\Molar} \ce{HNO3}:
        \answerblank[2.4cm]{2.60}
  \item pH of \SI{0.050}{\Molar} \ce{KOH}:
        \answerblank[2.4cm]{12.70}
  \item pH of \SI{0.0050}{\Molar} \ce{Ca(OH)2}: \workspace[2cm]
        \keyonly{\begin{apcanswer}$[\ce{OH-}] = 2(0.0050) = 0.010$;
        pOH $= 2.00$; pH $= \mathbf{12.00}$\end{apcanswer}}
  \item A student reports pH 12.00 for problem 3 by using
        $[\ce{OH-}] = 0.0050$. Do they get the right answer by the right
        method? \answerlines[2]{No --- they would get pOH $= 2.30$ and
        pH $= 11.70$. The correct answer requires doubling the hydroxide
        first, and the agreement of the final number here is
        coincidental.}
\end{enumerate}

\begin{exitticket}
Pure water at \SI{50}{\celsius} has pH 6.6. Is it acidic? Explain.
\keyonly{\begin{apcanswer}No, it is neutral. $K_w$ increases with
temperature, so both $[\ce{H+}]$ and $[\ce{OH-}]$ rise and the pH falls
below 7 --- but they remain \emph{equal}, which is what neutrality
means.\end{apcanswer}}
\end{exitticket}

% =====================================================================
\blocklesson{Weak Acid Equilibria}{CED 8.3 \textbullet\ Zumdahl \S14.5}

\begin{objectives}
  \item Calculate the pH of a weak acid solution.
  \item Compute and interpret percent ionization.
\end{objectives}

\reading{Zumdahl \S14.5}{710--719}

\begin{warmup}[3]
  \item pH of \SI{0.010}{\Molar} \ce{HCl}: \answerblank[2cm]{2.00}
  \item pH of \SI{0.010}{\Molar} \ce{Ba(OH)2}:
        \answerblank[2cm]{12.30}
  \item ICE stands for: \answerblank[5.8cm]{Initial, Change, Equilibrium}
\end{warmup}

\segment{INSTRUCTION A}{25}{A weak acid problem is a Unit 7 problem}

\notesection{$K_a$ and the ICE table}{8.3}
\practice{5}

\[ \ce{HA(aq) <=> H+(aq) + A^-(aq)} \qquad
   K_a = \frac{[\ce{H+}][\ce{A^-}]}{[\ce{HA}]} \]

Because $K_a$ is small, $x$ is usually negligible against $C_0$:
\[ K_a \approx \frac{x^2}{C_0} \quad\Longrightarrow\quad
   x = [\ce{H+}] = \sqrt{K_aC_0} \]

\begin{apctrap}
\textbf{The 5\% check is part of the answer.} Compute $x/C_0$ as a
percentage and compare with 5\%. An FRQ that asks you to justify the
approximation wants the number, not the assertion --- exactly as in Unit 7.
\end{apctrap}

\begin{workedexample}[Worked example 2: strong is not concentrated]
\SI{0.10}{\Molar} acetic acid, $K_a = 1.8\times10^{-5}$:
\[ x = \sqrt{(1.8\times10^{-5})(0.10)} = 1.34\times10^{-3}
   \quad\Longrightarrow\quad \mathrm{pH} = \mathbf{2.87} \]
Check: $1.34\times10^{-3}/0.10 = 1.3\%$ \checkmark

\textbf{Compare with \SI{0.10}{\Molar} \ce{HCl}, which has pH 1.00.} Same
concentration; the difference is entirely \emph{strength} --- the fraction
that ionizes. ``Strong'' and ``concentrated'' are different words describing
different things.
\end{workedexample}

\segment{APPLICATION}{20}{Weak acid calculations}

\begin{enumerate}[leftmargin=*,itemsep=0.4em]
  \item pH of \SI{0.250}{\Molar} \ce{HF} ($K_a = 7.2\times10^{-4}$),
        with the check. \workspace[2.4cm]
        \keyonly{\begin{apcanswer}$x = \sqrt{(7.2\times10^{-4})(0.250)} =
        1.34\times10^{-2}$; pH $= \mathbf{1.87}$.
        Check: $5.4\%$ --- marginally above 5\%, so strictly the quadratic
        should be used\end{apcanswer}}
  \item Percent ionization of \SI{0.100}{\Molar} acetic acid:
        \answerblank[2.4cm]{1.3\%}
  \item Diluting a weak acid raises the percent ionization but also raises
        the pH. Explain how both can be true.
        \answerlines[3]{$x \approx \sqrt{K_aC_0}$ falls only as the square
        root of concentration while $C_0$ falls linearly, so the
        \emph{fraction} ionized grows on dilution. But the absolute
        $[\ce{H+}]$ is still smaller, so the pH is higher.}
\end{enumerate}

\begin{exitticket}
Two solutions both have \SI{0.10}{\Molar} concentration: one \ce{HCl}, one
acetic acid. Which is more concentrated, and which is more acidic?
\keyonly{\begin{apcanswer}They are equally concentrated --- that is what
\SI{0.10}{\Molar} means. The \ce{HCl} is far more acidic, pH 1.00 against
2.87, because it ionizes completely while acetic acid ionizes only about
1.3\%.\end{apcanswer}}
\end{exitticket}

% =====================================================================
\blocklesson{Weak Bases and Polyprotic Acids}{CED 8.3 \textbullet\ Zumdahl \S14.6--14.7}

\begin{objectives}
  \item Calculate the pH of a weak base solution.
  \item Explain why the first ionization dominates for a polyprotic acid.
\end{objectives}

\reading{Zumdahl \S14.6--14.7}{719--730}

\begin{warmup}[3]
  \item $[\ce{H+}]$ for a weak acid: \answerblank[2.6cm]{$\sqrt{K_aC_0}$}
  \item The approximation threshold: \answerblank[1.6cm]{5\%}
  \item pH of \SI{0.10}{\Molar} acetic acid: \answerblank[2cm]{2.87}
\end{warmup}

\segment{INSTRUCTION A}{25}{The same table, one extra step}

\notesection{Weak bases and $K_b$}{8.3}
\practice{5}

\[ \ce{B(aq) + H2O(l) <=> BH+(aq) + OH^-(aq)} \qquad
   K_b = \frac{[\ce{BH+}][\ce{OH^-}]}{[\ce{B}]} \]

\begin{apctrap}
\textbf{For a base, the ICE $x$ is $[\ce{OH-}]$, not $[\ce{H+}]$.} Taking
$-\log x$ gives the \answerblank[1.6cm]{pOH}, and you must then subtract
from 14.00.

This is the single most-lost point in the unit. Reporting $-\log x$ as the
pH turns a strongly basic answer into a strongly acidic one --- a swing of
several units in the wrong direction, which any sanity check should catch.
\end{apctrap}

\begin{workedexample}[Worked example 3: ammonia]
\SI{0.100}{\Molar} \ce{NH3}, $K_b = 1.8\times10^{-5}$:
\[ x = \sqrt{(1.8\times10^{-5})(0.100)} = 1.34\times10^{-3} = [\ce{OH-}] \]
\[ \mathrm{pOH} = 2.87 \quad\Longrightarrow\quad
   \mathrm{pH} = 14.00 - 2.87 = \mathbf{11.13} \]
Note the pOH is numerically the same as acetic acid's pH --- the two have
equal constants. The final answers are wildly different, and the conversion
step is the only reason.
\end{workedexample}

\segment{INSTRUCTION B}{20}{Polyprotic acids}

\notesection{Why only the first ionization matters}{8.3}
\practice{5}

\ce{H3PO4} has $K_{a1} = 7.5\times10^{-3}$,
$K_{a2} = 6.2\times10^{-8}$, $K_{a3} = 4.8\times10^{-13}$.

Each successive constant is about $10^5$ times
\answerblank[2cm]{smaller}, because each proton must be pulled from an
increasingly \answerblank[2.4cm]{negative} anion. So the pH may be
calculated from \answerblank[2.4cm]{$K_{a1}$ alone}.

The exception is \ce{H2SO4}: its first ionization is
\answerblank[1.8cm]{strong} and its second
($K_{a2} = 1.2\times10^{-2}$) still contributes measurably.

\segment{APPLICATION}{20}{Bases and polyprotics}

\begin{enumerate}[leftmargin=*,itemsep=0.4em]
  \item pH of \SI{0.150}{\Molar} \ce{NH3}. \workspace[2.4cm]
        \keyonly{\begin{apcanswer}$x = \sqrt{(1.8\times10^{-5})(0.150)} =
        1.64\times10^{-3}$; pOH $= 2.78$;
        pH $= \mathbf{11.22}$\end{apcanswer}}
  \item A student reports pH 2.78 for that problem. Identify the error.
        \answerlines[2]{They reported the pOH as the pH. For a base the ICE
        $x$ is $[\ce{OH-}]$, so $-\log x$ is the pOH and the pH is
        $14.00 - 2.78 = 11.22$.}
  \item Justify calculating \ce{H3PO4}'s pH from $K_{a1}$ alone.
        \answerlines[2]{$K_{a2}$ is roughly $10^5$ times smaller, so the
        second ionization contributes a negligible quantity of \ce{H+}
        beside the first.}
\end{enumerate}

\begin{exitticket}
Given $K_a$ for a weak acid, how do you find $K_b$ for its conjugate base,
and why does the relationship hold?
\keyonly{\begin{apcanswer}$K_b = K_w/K_a$. Adding the acid ionization to the
conjugate base hydrolysis gives the autoionization of water, and adding
reactions multiplies their equilibrium constants --- so
$K_aK_b = K_w$.\end{apcanswer}}
\end{exitticket}

% =====================================================================
\blocklesson{Structure, p$K_a$, and Salts}{CED 8.6--8.7 \textbullet\ Zumdahl \S14.8--14.9}

\begin{objectives}
  \item Explain acid strength trends from molecular structure.
  \item Predict whether a salt solution is acidic, basic, or neutral.
\end{objectives}

\reading{Zumdahl \S14.8--14.9}{730--740}

\begin{warmup}[3]
  \item For a weak base, ICE $x$ is: \answerblank[2.2cm]{$[\ce{OH-}]$}
  \item pH of \SI{0.100}{\Molar} \ce{NH3}: \answerblank[2cm]{11.13}
  \item $K_b$ for acetate: \answerblank[3.2cm]{$5.6\times10^{-10}$}
\end{warmup}

\segment{INSTRUCTION A}{25}{p$K_a$ and structure}

\notesection{p$K_a$}{8.7}
\practice{5}

\[ \mathrm{p}K_a = -\log K_a \]
A \emph{lower} p$K_a$ means a \answerblank[2.2cm]{stronger} acid. Acetic
acid's $K_a = 1.8\times10^{-5}$ gives p$K_a = \answerblank[1.6cm]{4.74}$.

\notesection{Structure and acid strength}{8.6}
\practice{6}

\begin{itemize}[leftmargin=*,itemsep=0.3em]
  \item \textbf{Binary acids down a group:} \ce{HF} $<$ \ce{HCl} $<$
        \ce{HBr} $<$ \ce{HI}. The controlling factor is
        \answerblank[2.8cm]{bond strength} --- the H--X bond weakens down
        the group.
  \item \textbf{Oxyacids:} \ce{HClO} $<$ \ce{HClO2} $<$ \ce{HClO3} $<$
        \ce{HClO4}. More electronegative oxygens withdraw electron density
        from the O--H bond \emph{and} stabilize the conjugate base.
\end{itemize}

\begin{apctrap}
\ce{HF} is the \emph{weakest} hydrohalic acid despite fluorine being the
most electronegative element. Holding the proton tightly in a very strong
H--F bond matters more than pulling electron density away from it. Do not
reach for electronegativity when the trend is down a group.
\end{apctrap}

\segment{INSTRUCTION B}{20}{Salt solutions}

\notesection{Trace each ion to its parent}{8.7}
\practice{6}

\begin{center}
\begin{tabular}{@{}lllll@{}}
\hline
\textbf{Cation from} & \textbf{Anion from} & \textbf{Solution} &
  \textbf{Example} \\
\hline
strong base & strong acid & \textbf{neutral} & \ce{NaCl} \\
strong base & weak acid & \textbf{basic} & \ce{NaC2H3O2} \\
weak base & strong acid & \textbf{acidic} & \ce{NH4Cl} \\
\hline
\end{tabular}
\end{center}

The logic is entirely conjugate-pair reasoning: an anion from a
\emph{weak} acid is a meaningful base, while one from a \emph{strong} acid
is not.

\segment{APPLICATION}{20}{Structure and salts}

\begin{enumerate}[leftmargin=*,itemsep=0.4em]
  \item Classify and name the responsible ion: \ce{KNO3}, \ce{NH4NO3},
        \ce{NaCN}. \answerlines[2]{\ce{KNO3} neutral --- neither ion
        reacts. \ce{NH4NO3} acidic --- \ce{NH4+}. \ce{NaCN} basic ---
        \ce{CN-}.}
  \item Calculate the pH of \SI{0.30}{\Molar} \ce{NaF}
        ($K_a(\ce{HF}) = 7.2\times10^{-4}$). \workspace[2.8cm]
        \keyonly{\begin{apcanswer}$K_b = 1.0\times10^{-14}/7.2\times10^{-4}
        = 1.4\times10^{-11}$;
        $x = \sqrt{(1.4\times10^{-11})(0.30)} = 2.0\times10^{-6}$;
        pOH $= 5.69$; pH $= \mathbf{8.31}$\end{apcanswer}}
  \item Explain why \ce{NaCl} solutions are neutral but \ce{NaCN}
        solutions are basic. \answerlines[3]{Both contain \ce{Na+}, from a
        strong base, which does not react. \ce{Cl-} comes from the strong
        acid \ce{HCl} and is a negligible base. \ce{CN-} comes from the
        very weak acid \ce{HCN}, so it readily takes a proton from water
        and generates \ce{OH-}.}
\end{enumerate}

\begin{exitticket}
Rank \ce{HF}, \ce{HCl}, \ce{HBr}, \ce{HI} by increasing strength and name
the controlling factor.
\keyonly{\begin{apcanswer}\ce{HF} $<$ \ce{HCl} $<$ \ce{HBr} $<$ \ce{HI}.
The H--X bond weakens down the group, so the proton is released more
readily; bond strength outweighs electronegativity
here.\end{apcanswer}}
\end{exitticket}

% =====================================================================
\blocklesson{Acid--Base Reactions and Buffers}{CED 8.4 \textbullet\ Zumdahl \S15.1--15.2}

\begin{objectives}
  \item Predict the products of an acid--base reaction.
  \item Recognize a buffer and explain the common-ion effect.
\end{objectives}

\reading{Zumdahl \S15.1--15.2}{756--765}

\begin{warmup}[3]
  \item p$K_a$ of acetic acid: \answerblank[1.8cm]{4.74}
  \item \ce{NH4Cl} solutions are: \answerblank[2cm]{acidic}
  \item Adding a product shifts an equilibrium:
        \answerblank[1.8cm]{left}
\end{warmup}

\segment{INSTRUCTION A}{25}{The common-ion effect}

\notesection{Le Ch\^atelier applied to an acid}{8.4}
\practice{6}

\[ \ce{HF(aq) <=> H+(aq) + F^-(aq)} \]
Dissolve \ce{NaF} in this solution and you have added a
\answerblank[2cm]{product}, so the equilibrium shifts
\answerblank[1.6cm]{left} and the acid ionizes \emph{less}.

The algebra collapses to a ratio:
\[ [\ce{H+}] = K_a\frac{[\ce{HA}]}{[\ce{A^-}]} \]
No square root, no quadratic --- when both members of a conjugate pair are
present in comparable amounts, $x$ is negligible against both.

\segment{INSTRUCTION B}{20}{What makes a buffer}

\notesection{Recognizing one}{8.4}
\practice{6}

A \term{buffer} contains significant amounts of \textbf{both members of a
conjugate pair}. Three ways it appears:

\begin{enumerate}[leftmargin=*,itemsep=0.25em]
  \item Stated outright --- \ce{HC2H3O2} plus \ce{NaC2H3O2}.
  \item Made by \answerblank[4.6cm]{partial neutralization} of a weak acid
        with strong base.
  \item Hiding at the halfway point of a titration.
\end{enumerate}

\begin{apctrap}
\ce{HCl} with \ce{NaCl} is \textbf{not} a buffer. \ce{Cl-} has no affinity
for protons, so nothing absorbs added acid. A buffer needs a \emph{weak}
conjugate pair --- reversibility is the whole point.
\end{apctrap}

\segment{APPLICATION}{20}{Recognition and prediction}

\begin{enumerate}[leftmargin=*,itemsep=0.35em]
  \item Buffer or not, with a reason: \ce{NH3} $+$ \ce{NH4Cl}
        \answerblank[8.3cm]{yes --- weak base and its conjugate acid}
  \item \ce{NaOH} $+$ \ce{NaCl}
        \answerblank[8cm]{no --- no weak conjugate pair}
  \item Find $[\ce{H+}]$ and pH for a solution \SI{0.10}{\Molar} in
        \ce{HF} and \SI{0.30}{\Molar} in \ce{NaF}. \workspace[2.2cm]
        \keyonly{\begin{apcanswer}$[\ce{H+}] = (7.2\times10^{-4})(0.10/0.30)
        = 2.4\times10^{-4}$; pH $= \mathbf{3.62}$\end{apcanswer}}
\end{enumerate}

\begin{exitticket}
Explain why adding \ce{NaF} to an \ce{HF} solution raises its pH.
\keyonly{\begin{apcanswer}\ce{F-} is a product of the \ce{HF} ionization, so
adding it shifts that equilibrium to the left and suppresses the production
of \ce{H+}. The hydrolysis of \ce{F-} contributes as well, but the dominant
effect is the common-ion shift.\end{apcanswer}}
\end{exitticket}

% =====================================================================
\blocklesson{Henderson--Hasselbalch}{CED 8.8--8.9 \textbullet\ Zumdahl \S15.2}

\begin{objectives}
  \item Apply the Henderson--Hasselbalch equation.
  \item Explain why the ratio, not the concentration, sets the pH.
\end{objectives}

\reading{Zumdahl \S15.2}{763--767}

\begin{warmup}[3]
  \item $[\ce{H+}] = K_a \times$
        \answerblank[3.4cm]{$[\ce{HA}]/[\ce{A^-}]$}
  \item A buffer contains: \answerblank[8.3cm]{both members of a conjugate
        pair}
  \item p$K_a$ of \ce{NH4+}: \answerblank[1.8cm]{9.25}
\end{warmup}

\segment{INSTRUCTION A}{25}{Taking the logarithm}

\notesection{The equation}{8.9}
\practice{5}

\[ \boxed{\;\mathrm{pH} = \mathrm{p}K_a
   + \log\frac{[\text{base}]}{[\text{acid}]}\;} \]

\begin{itemize}[leftmargin=*,itemsep=0.3em]
  \item Equal amounts $\Rightarrow$ log term is
        \answerblank[1.4cm]{0}, so pH $=$
        \answerblank[1.8cm]{p$K_a$}.
  \item More base than acid $\Rightarrow$ pH
        \answerblank[3cm]{above p$K_a$}.
  \item More acid than base $\Rightarrow$ pH
        \answerblank[3cm]{below p$K_a$}.
\end{itemize}

\begin{apctrap}
\textbf{Base over acid, and the base is the anion.} Inverting the ratio
flips the sign of the correction, landing the pH the same distance on the
wrong side of p$K_a$. Before dividing, name the two species aloud:
\ce{C2H3O2^-} is the base, \ce{HC2H3O2} is the acid.
\end{apctrap}

\segment{INSTRUCTION B}{20}{Why the ratio is what matters}

\notesection{Dilution barely moves a buffer's pH}{8.8}
\practice{5}

Only the \emph{ratio} appears, so:

\begin{itemize}[leftmargin=*,itemsep=0.3em]
  \item Diluting a buffer leaves the pH essentially
        \answerblank[2.4cm]{unchanged} --- both concentrations fall by the
        same factor. Contrast a weak acid alone, where dilution
        \emph{does} raise the pH.
  \item You may use \answerblank[1.8cm]{moles} instead of molarity, since
        the shared volume cancels. This saves a step on nearly every
        titration problem.
\end{itemize}

\begin{workedexample}[Worked example 4: two buffers, one pH]
\SI{0.25}{\Molar} \ce{HC2H3O2} with \SI{0.15}{\Molar} \ce{NaC2H3O2}:
\[ \mathrm{pH} = 4.74 + \log(0.15/0.25) = 4.74 - 0.22 = \mathbf{4.52} \]

\SI{0.050}{\Molar} with \SI{0.030}{\Molar} gives the same ratio $0.60$ and
therefore the \textbf{same pH, 4.52} --- at one fifth the concentration. The
pH is identical; the \emph{capacity} is not.
\end{workedexample}

\segment{APPLICATION}{20}{Buffer calculations}

\begin{enumerate}[leftmargin=*,itemsep=0.4em]
  \item pH of \SI{0.20}{\Molar} \ce{NH3} with \SI{0.30}{\Molar}
        \ce{NH4Cl}. \workspace[2cm]
        \keyonly{\begin{apcanswer}pH $= 9.25 + \log(0.20/0.30) =
        9.25 - 0.18 = \mathbf{9.07}$\end{apcanswer}}
  \item What ratio $[\ce{A^-}]/[\ce{HA}]$ gives pH 5.00 for acetic acid?
        \answerblank[2.4cm]{1.8}
  \item Two acetate buffers: X is \SI{1.0}{\Molar} in each component,
        Y is \SI{0.010}{\Molar} in each. Compare their pH values and their
        usefulness. \answerlines[3]{Both have pH 4.74 --- the ratio is 1 in
        each. X is far more useful: it holds 100 times as much of each
        component, so it can absorb much more added acid or base before the
        ratio moves appreciably.}
\end{enumerate}

\begin{exitticket}
Why does diluting a buffer barely change its pH, while diluting a weak acid
alone raises it noticeably?
\keyonly{\begin{apcanswer}A buffer's pH depends on the \emph{ratio} of two
concentrations, and dilution divides both by the same factor, leaving it
unchanged. A weak acid alone has $[\ce{H+}] = \sqrt{K_aC_0}$, which depends
on $C_0$ itself, so lowering $C_0$ lowers $[\ce{H+}]$.\end{apcanswer}}
\end{exitticket}

% =====================================================================
\blocklesson{Buffer Action and Capacity}{CED 8.10 \textbullet\ Zumdahl \S15.2--15.3}

\begin{objectives}
  \item Calculate a buffer's pH after strong acid or base is added.
  \item Compare buffering capacity and design a buffer.
\end{objectives}

\reading{Zumdahl \S15.2--15.3}{761--770}

\begin{warmup}[3]
  \item Henderson--Hasselbalch: pH $=$
        \answerblank[5.4cm]{p$K_a + \log$(base/acid)}
  \item pH when the ratio is 1: \answerblank[2cm]{p$K_a$}
  \item Diluting a buffer changes the pH how?
        \answerblank[3cm]{hardly at all}
\end{warmup}

\segment{INSTRUCTION A}{25}{Stoichiometry first, then equilibrium}

\notesection{The two-step method}{8.10}
\practice{5}

\begin{enumerate}[leftmargin=*,itemsep=0.3em]
  \item \textbf{Stoichiometry.} The strong reagent reacts
        \answerblank[2.6cm]{completely}. Added \ce{OH-} converts \ce{HA} to
        \ce{A-}; added \ce{H+} does the reverse. Work in moles. No $K$
        appears.
  \item \textbf{Equilibrium.} Put the new amounts into
        Henderson--Hasselbalch.
\end{enumerate}

\begin{apctrap}
\textbf{Never put added strong base into an ICE table.} Strong base does not
sit at equilibrium with a weak acid --- it consumes it. Mixing the two steps
is the single most reliable way to lose an entire buffer question.
\end{apctrap}

\begin{workedexample}[Worked example 5: a buffer absorbing a base]
\SI{1.00}{\litre} of a buffer is \SI{0.50}{\Molar} in \ce{HC2H3O2} and
\SI{0.50}{\Molar} in \ce{NaC2H3O2}. Add \SI{0.010}{\mole} \ce{NaOH}.

\textbf{Start:} ratio 1, so pH $= 4.74$.

\textbf{Step 1:} $\ce{OH^- + HC2H3O2 -> C2H3O2^- + H2O}$ leaves
$0.49$ mol acid and $0.51$ mol base.

\textbf{Step 2:} pH $= 4.74 + \log(0.51/0.49) = \mathbf{4.76}$.

\textbf{For contrast,} the same \SI{0.010}{\mole} of \ce{NaOH} in
\SI{1.00}{\litre} of pure water gives pH $\mathbf{12.00}$. The buffer moved
\textbf{0.02}; water moved \textbf{5.00}.
\end{workedexample}

\segment{INSTRUCTION B}{20}{Capacity and design}

\notesection{Why equal amounts buffer best}{8.10}
\practice{6}

\term{Buffering capacity} depends on two things: how \emph{much} of each
component is present, and how \emph{balanced} the ratio is.

\begin{center}
\fbox{\parbox{0.86\linewidth}{\centering
Optimal buffering when $[\ce{HA}] = [\ce{A^-}]$, i.e.\ pH $=$ p$K_a$. \\
Useful range: p$K_a \pm 1$.}}
\end{center}

To \textbf{design} a buffer, choose a weak acid whose p$K_a$ is as close as
possible to the target pH, then set the ratio with
Henderson--Hasselbalch.

\segment{APPLICATION}{20}{Action and design}

\begin{enumerate}[leftmargin=*,itemsep=0.4em]
  \item Add \SI{0.010}{\mole} \ce{HCl} to the buffer in worked example 5
        instead. New pH? \answerblank[2.4cm]{4.72}
  \item Design a pH 4.00 buffer given p$K_a$ values \ce{HNO2} 3.40,
        \ce{HC2H3O2} 4.74, \ce{HOCl} 7.46. \workspace[2.4cm]
        \keyonly{\begin{apcanswer}\ce{HNO2} is closest (0.60 away).
        $4.00 = 3.40 + \log r$, so $r = 10^{0.60} = \mathbf{4.0}$ --- four
        times as much nitrite as nitrous acid.\end{apcanswer}}
  \item Why would \ce{HOCl} be a poor choice there?
        \answerlines[2]{Its p$K_a$ is 3.5 units away, so the required ratio
        is about $10^{-3.5}$ --- almost no conjugate base present, and
        therefore essentially no capacity against added acid.}
\end{enumerate}

\begin{exitticket}
State the two steps for a buffer receiving strong base, and say what makes
each different.
\keyonly{\begin{apcanswer}First a stoichiometry step: the strong base reacts
to completion, converting weak acid into conjugate base, worked in moles
with no equilibrium constant. Then an equilibrium step: substitute the new
amounts into Henderson--Hasselbalch.\end{apcanswer}}
\end{exitticket}

% =====================================================================
\blocklesson{Titrations, Indicators, and pH Effects}{CED 8.5, 8.11 \textbullet\ Zumdahl \S15.4--15.5, \S16.1}

\begin{objectives}
  \item Analyse a titration curve region by region.
  \item Select an indicator and predict pH effects on solubility.
\end{objectives}

\reading{Zumdahl \S15.4--15.5}{770--792} \reading{Zumdahl \S16.1}{812--813}

\begin{warmup}[3]
  \item Buffer method step 1: \answerblank[7.3cm]{stoichiometry to
        completion}
  \item Optimal buffering ratio: \answerblank[1.6cm]{1}
  \item Useful buffer range: \answerblank[2.4cm]{p$K_a \pm 1$}
\end{warmup}

\segment{INSTRUCTION A}{25}{Four regions, four methods}

\notesection{Weak acid titrated with strong base}{8.5}
\practice{5}

\begin{center}
\renewcommand{\arraystretch}{1.2}
\begin{tabular}{@{}llp{5.4cm}@{}}
\hline
\textbf{Region} & \textbf{Major species} & \textbf{Method} \\
\hline
before any base & \ce{HA} only & weak-acid ICE \\
before equivalence & \ce{HA} $+$ \ce{A-} & \textbf{buffer:} H--H \\
at equivalence & \ce{A-} only & weak-\emph{base} ICE, $K_b = K_w/K_a$ \\
after equivalence & excess \ce{OH-} & strong base alone \\
\hline
\end{tabular}
\end{center}

\begin{apcnote}[Two landmarks worth a point each]
\textbf{Halfway point:} $[\ce{HA}] = [\ce{A^-}]$, so pH $=$ p$K_a$. This is
how $K_a$ is measured from a curve.

\textbf{Equivalence point:} moles of base $=$ moles of acid. That is a
\emph{stoichiometric} definition and says nothing about the pH being 7.
\end{apcnote}

\begin{apctrap}
\textbf{A weak-acid equivalence point is not pH 7.} All the acid has become
its conjugate base, which hydrolyses. Weak acid with strong base gives an
equivalence pH \answerblank[2cm]{above 7}; weak base with strong acid gives
one \answerblank[2cm]{below 7}. Only strong-with-strong lands on 7.00.
\end{apctrap}

\begin{workedexample}[Worked example 6: acetic acid titration]
\SI{50.0}{\milli\litre} of \SI{0.10}{\Molar} \ce{HC2H3O2}
($= \SI{5.0}{\milli\mole}$) with \SI{0.10}{\Molar} \ce{NaOH}; equivalence
at \SI{50.0}{\milli\litre}.

\begin{center}
\begin{tabular}{@{}rlr@{}}
\hline
\textbf{mL} & \textbf{Region} & \textbf{pH} \\
\hline
0 & weak acid & 2.87 \\
25.0 & \textbf{halfway} --- pH $=$ p$K_a$ & \textbf{4.74} \\
50.0 & \textbf{equivalence} --- \ce{A-} only & \textbf{8.72} \\
60.0 & excess base & 11.96 \\
\hline
\end{tabular}
\end{center}

At equivalence, $[\ce{C2H3O2^-}] = 5.0/100.0 = 0.050$~M and
$K_b = 5.6\times10^{-10}$, giving $[\ce{OH-}] = 5.3\times10^{-6}$,
pOH $= 5.28$, pH $= 8.72$ --- basic, as it must be.
\end{workedexample}

\segment{INSTRUCTION B}{20}{Indicators and pH effects}

\notesection{Choosing an indicator}{8.5}
\practice{6}

An indicator is itself a weak acid whose two forms differ in colour. It
changes over roughly \answerblank[2.6cm]{p$K_a \pm 1$}, so choose one whose
range brackets the \answerblank[3.4cm]{equivalence pH}.

The \term{equivalence point} is fixed by stoichiometry; the
\term{end point} is where the indicator changes. A good choice makes them
nearly coincide.

\notesection{pH and solubility}{8.11}
\practice{6}

\begin{apcnote}[Qualitative only]
CED topic 8.11 carries an \textbf{exclusion statement}: computations of
solubility as a function of pH are never assessed. Predict the direction and
justify with Le Ch\^atelier --- do not calculate.
\end{apcnote}

A salt is more soluble in acid when its anion is a
\answerblank[3.6cm]{meaningful base} --- that is, when its conjugate acid is
weak. \ce{CaCO3} dissolves in acid because \ce{CO3^2-} is consumed;
\ce{AgCl} does not, because \ce{Cl-} comes from a strong acid and does not
react.

\segment{APPLICATION}{20}{Curves, indicators, and solubility}

\begin{enumerate}[leftmargin=*,itemsep=0.4em]
  \item A curve for a monoprotic acid shows equivalence at pH 9.1. Is the
        acid strong or weak? Justify. \answerlines[2]{Weak. Strong with
        strong gives exactly 7.00; a basic equivalence point means the
        anion left behind hydrolyses, which happens only for the conjugate
        base of a weak acid.}
  \item Choose an indicator for that titration: methyl red (4.4--6.2) or
        phenolphthalein (8.2--10.0)?
        \answerblank[3.4cm]{phenolphthalein}
  \item \ce{BaSO4} and \ce{CaCO3} are each treated with dilute acid.
        Predict and justify. \answerlines[3]{\ce{CaCO3} becomes more
        soluble --- \ce{CO3^2-} is the conjugate base of the weak acid
        \ce{H2CO3}, so \ce{H+} removes it and the dissolution equilibrium
        shifts right. \ce{BaSO4} is essentially unaffected, since
        \ce{SO4^2-} comes from a strong acid and is a negligible base.}
\end{enumerate}

\begin{exitticket}
Give the single question to ask before starting any calculation in this
unit, and why it settles the method.
\keyonly{\begin{apcanswer}``What are the major species after any strong
reagent has reacted completely?'' The answer places you in one of four
cases --- strong acid or base, weak acid or base alone, a buffer, or excess
strong reagent --- and each has its own method. Choosing a method before
identifying the species is how students end up applying
Henderson--Hasselbalch at an equivalence point, where no buffer
exists.\end{apcanswer}}
\end{exitticket}

\end{document}
